\section{Related Work}\label{sec:related}


\mypara{SMT-Based Verification}
%
SMT-solvers have been extensively
used to automate program verification
via Floyd-Hoare logics~\cite{Nelson81}.
%
\toolfont{Leon} introduces an
SMT-based algorithm that is
complete for catamorphisms
(folds) over ADTs~\cite{Suter2010}
and a semi-decision procedure that 
is guaranteed to find satisfying assignments 
(models) for queries over arbitrary recursive
functions, if they exist~\cite{Suter2011}.
%
Our work is inspired by \dafny's Verified
Calculations~\citep{LeinoPolikarpova16} 
but differs in
%
(1)~our use of reflection instead of axiomatization,
%
(2)~our use of refinements to compose proofs, and
%
(3)~our use of \pbesym to automate reasoning
about user-defined functions.
%
\dafny (and \fstar~\citep{fstar})
encode user-functions as axioms and use
a fixed fuel to instantiate functions
upto some fixed unfolding
depth~\cite{Amin2014ComputingWA}.
%
While the fuel-based approach is incomplete,
even for equational or calculational reasoning,
it may, although rare in practice, quickly
time out after a fixed, small number of
instantiations rather than perform an
exhaustive proof search like \pbesym.
%
Nevertheless, \pbesym demonstrates that it
is possible to develop complete and practical
algorithms for reasoning about user-defined
functions.

\mypara{Proving Equational Properties}
%
Several authors have proposed tools for proving
(equational) properties of (functional) programs.
%
Systems of~\citet{sousa16} and~\citet{KobayashiRelational15}
extend classical safety verification algorithms,
respectively based on Floyd-Hoare logic and refinement types,
to the setting of relational or $k$-safety properties
that are assertions over $k$-traces of a program.
%
Thus, these methods can automatically prove that
certain functions are associative, commutative \etc
but are restricted to first-order properties and
are not programmer-extensible.
%
Zeno~\citep{ZENO} generates proofs by term
rewriting and Halo~\citep{HALO} uses an axiomatic
encoding to verify contracts.
%
Both the above are automatic, but unpredictable and not
programmer-extensible, hence, have been limited to
far simpler properties than the ones checked here.
%
HERMIT~\citep{Farmer15} proves equalities by rewriting
the GHC core language, guided by user specified scripts.
%
Our proofs are Haskell programs, SMT solvers
automate reasoning, and importantly, we
connect the validity of proofs with the
semantics of the programs.


% Dafny, and the related \fstar~\citep{fstar}
% which like \toolname, uses types to compose
% proofs, offer more automation by translating
% recursive functions to SMT axioms.
% However, unlike reflection, this axiomatic
% approach renders typechecking and verification
% undecidable (in theory) and leads to
% unpredictability and divergence
% (in practice)~\citep{Leino16}.

% \RJ{PBE brings completeness for the first time yada}
%
% \RJ{More stuff from "butterfly effect" text? (see scratch.tex)}
%
%\NV{CHECL Relational-F*, Barthe et al, from POPL 2014, and EasyCrypt}

%% In a work more closely related to
%% ours, \fstar uses refinement types
%% for program verification supporting
%% expressiveness of fully dependent types.
%% %
%% As in Dafny, \fstar directly translates
%% recursive functions to axioms in the logic
%% thus suffers from the ``butterfly effect''
%% and allows the user to explicitly write SMT tactics to control it.

%% Leino \etal~\citep{Leino16}
%% name this problem as the ``butterfly
%% effect'', in which minor modifications
%% to the program source cause significant
%% instabilities in verification and propose
%% trigger selection strategies to address it.
%% %
%% We avoid the ``butterfly effect'' by not
%% directly axiomatizing functions into logic.
%% Instead the information about the function's
%% body is exactly captured in function's result
%% type and user needs to explicitly invoke the function to push
%% the function's definition information into the logic.


% mature proof assistants like \coq and \isabelle
% could be adapted to the refinement setting.
%
% that, unlike our approach, rely on tiny trusted computing base.
%%%enable large scale proof engineering
%%%that is well beyond what is currently
% beyond what is currently feasible with our approach.
%
% (Indeed, it is likely that our proof search procedure
% can, like @omega@, be encoded as a tactic within \coq.)
% which allow for arbitrary expressiveness of the type system
% in the cost of automatic verification.
%
%% The syntax of \libname's operators is inspired by
%% Equational Reasoning in Agda~\citep{agda}.
%% Here we extended these equational operators
%% to support linear arithmetic and, for example, prove
%% properties of Ackermann function.
%% %
%% Unlike Adga, proof term are explicit in \libname,
%% we do not use heuristics to infer proofs.


%%\mypara{Dependent Types for Non-Terminating Programs}
%%%
%%Zombie~\citep{Zombie, Sjoberg2015} integrates
%%dependent types in non terminating programs
%%and supports automatic reasoning for equality.
%%%
%%Vazou \etal have previously~\citep{Vazou14} shown
%%how Liquid Types can be used to check
%%non-terminating programs.
%%%
%%Reflection makes \toolname at least as
%%expressive as Zombie, \emph{without}
%%having to axiomatize the theory of
%%equality within the type system.
%%%
%%Consequently, in contrast to Zombie,
%%SMT based reflection lets \toolname
%%verify higher-order specifications
%%like @foldr_fusion@.

% which lets us use SMT automation
% to verify deep specifications of
% non-trivial programs.
%

\mypara{Dependent Types in Programming}
%
Integration of dependent types into Haskell
has been a long standing goal~\citep{EisenbergS14}
that dates back to Cayenne~\citep{cayenne},
a Haskell-like, fully dependent type
language with undecidable type checking.
%
% In a recent line of work~\citet{EisenbergS14}
% aim to allow fully dependent
% programming within Haskell, by making
% ``type-level programming ... at least as
%   expressive as term-level programming''.
%
Our approach differs significantly in that
reflection and \pbesym use SMT-solvers
to drastically simplify proofs over decidable
theories.
%
Zombie~\citep{Sjoberg2015} investigates
the design of a dependently typed language where
SMT-style congruence closure is used to reason
about the equality of terms.
%
However, Zombie explicitly eschews type-level
computation, as the authors write
% note that type level computation, or to be precise
``equalities that follow from $\beta$-reduction''
are ``incompatible with congruence closure''.
%
Due to this incompleteness, the programmer must
use explicit @join@ terms to indicate where
normalization should be triggered, even
so, equality checking is based on fuel,
hence, is incomplete.
%
% In contrast, reflection and \pbesym show
% how congruence closure and computation can
% be very profitably combined, freeing the
% programmer from having to provide
% explicit @join@ operations.
% \RJ{THE ABOVE IS TOO LONG}


% Second, refinements are completely erased at run-time.
% That is, while both systems automatically lift Haskell
% code to either uninterpreted logical functions
% or type families, with refinements, the logical
% functions are not accessible at run-time and
% promotion cannot affect the semantics of
% the program.
%
% As an advantage (resp. disadvantage), refinements
% cannot degrade (resp. optimize)
% the performance of programs.

\mypara{Theorem Provers}
%
Reflection shows how to retrofit deep
specification and verification in the
style of \agda~\citep{agda},
\coq~\citep{coq-book} and
\isabelle~\cite{NPW2002} into existing
languages via refinement typing and
\pbesym shows how type-level computation
can be made compatible with SMT solvers'
native theory reasoning yielding a
powerful new way to automate
proofs~(\S~\ref{sec:overview:lists}).
%
An extensive comparison~\cite{Vazou17}
between our approach and mature theorem 
provers like \coq, \agda, and \isabelle
reveals that these provers 
have two clear advantages
over our approach:
%
they emit certificates,
so they rely on a small
trusted computing base,
and they have decades-worth
of tactics, libraries,
and proof scripts that
enable large scale proof
engineering.
%
Some tactics even enable
embedding of SMT-based
proof search heuristics,
\eg \toolfont{Sledgehammer}~\cite{sledgehammer},
that is widely used in
\isabelle.
However, this search
does not have the
completeness
guarantees of \pbesym.
%
The issue of extracting checkable
certificates from SMT solvers is well
understood~\cite{Necula97,chen-pldi-2010}
and easy to extend to our setting.
%
However, the question of extending SMT-based
verifiers with tactics and scriptable proof
search, and more generally, incorporating
\emph{interactivity} in the style of
proof-assistants, perhaps enhanced by
proof-completion hints, remains an
interesting direction for future work.

% However, the question of integrating
% tactics and scriptable proof search
% within SMT-based verifiers remains an
% interesting direction for future work.



%% Compared to these systems, our proofs are
%% expressed as Haskell programs and do not
%% require the user to learn a different
%% tactic languages.
%% %
%% Moreover, our system is more general
%% as it allows for both equational
%% and linear arithmetic proofs.
%% %
%% On the other hand, \libname requires
%% explicit proofs and does not currently
%% support any automatic heuristics.

%%% CUT \mypara{Deterministic Parallelism}
%%% CUT %
%%% CUT Deterministic parallelism has plenty of theory but relatively few practical
%%% CUT implementations.  Early discoveries were based on limited producer-consumer
%%% CUT communication, such as single-assignment variables \cite{Tesler-1968,IStructures}, Kahn
%%% CUT process networks~\cite{kahn-1974}, and synchronous dataflow~\cite{lee-sdf}.
%%% CUT Other models use synchronous updates to shared state, as in
%%% CUT Esterel~\cite{synchronous-overview} or PRAM.  Finally, work on type systems for
%%% CUT permissions management \cite{permission-types,habanero-java-permissions},
%%% CUT supports the development of {\em non-interfering} parallel programs that access
%%% CUT disjoint subsets of the heap in parallel.  Parallel functional programming is
%%% CUT also non-interfering~\cite{manticore,multicore-ghc}.
%%% CUT %
%%% CUT Irrespective of which theory is used to support deterministic parallel
%%% CUT programming, practical implementations such as Cilk~\cite{cilk} or Intel
%%% CUT CnC~\cite{cnc} are limited by host languages with type systems insufficient to
%%% CUT limit side effects, much less prove associativity.  Conversely, dependently
%%% CUT typed languages like Agda and Idris do not have parallel programming APIs and
%%% CUT runtime systems.

% synchronous languages such as Esterel
